\section{Investiga}

\subsection{L'Escola d'Atenes, de Rafael}

[draft]

La Escuela de Atenas de Rafael, pintada hacia 1509–1511, es básicamente un meme intelectual del siglo XVI con presupuesto papal.

La idea era representar a los grandes cerebros de la Antigüedad reunidos en una especie de convención imposible: Platón, Aristóteles, Sócrates, Pitágoras, Euclides, Ptolomeo… todos juntos, siglos de diferencia aparte, como si hubieran quedado para discutir después de comer.

Y Rafael tuvo cierta coña. A varios filósofos les puso caras de gente de su época: Platón parece Leonardo da Vinci, Euclides probablemente Bramante y, como pulla particularmente elegante, pintó a Heráclito con la cara de Miguel Ángel, solo y bastante enfurruñado.

Además, Rafael se metió a sí mismo en una esquina mirando al espectador. Básicamente: «sí, yo también estaba aquí».

¿Por qué Platón apunta al cielo y Aristoteles extiende su mano hacia la tierra?

\subsection{Sara García Alonso}

[draft] sara alonso

Sara García Alonso es biotecnóloga, investigadora contra el cáncer en el CNIO y astronauta de reserva de la Agencia Espacial Europea. Y sí, académicamente es un pepino: doctora cum laude en Biología Molecular del Cáncer y seleccionada en 2022 entre más de 22.500 candidatos para formar parte de la nueva promoción de astronautas de la ESA.

Pero quizá lo más interesante de Sara no sea acumular medallas. Insiste mucho en algo bastante menos espectacular: la ciencia y el espacio son trabajos de equipo. Ella misma cuenta que aprender a colaborar le cambió la forma de trabajar y que tanto en un laboratorio como en una tripulación puedes llegar más lejos con otros que por tu cuenta.

Además, divulga bastante sobre cómo es el día a día real de un laboratorio: experimentos, investigación contra el cáncer, formación, errores, aprendizaje… No solo la foto bonita con el traje de astronauta. El CNIO incluso ha publicado vídeos suyos explicando cómo trabaja una investigadora del cáncer. Una científica bastante buena para tener en el radar.

Si te interesa puedes seguir a Sara Alonso, es muy activa en las redes y ha escrito un libro y ha dada muchas entrevistas.
